% Full proof: root, boundary_d, and limits_ab source reviews passed.
\subsection{No proper analytic filling of the actual punctured family}

\begin{theorem}[B.4.33; Proved; nonexistence after every finite parameter cover]
\label{push:B433}
The actual punctured six-sphere family has no proper
Hausdorff complex-analytic filling after any finite
cover of its degeneration parameter. In particular,
there is no positive integer $N$ for which the cover
$w=r^N$ has a proper filling with smooth total space
and reduced normal-crossings central fibre.

The assertion permits an arbitrary singular central
space, arbitrary additional components, and arbitrary
central multiplicities. It does not impose a
K\"ahler, projective, or logarithmic hypothesis on
the putative filling. Its proof derives the specific
logarithmic torsor needed for the contradiction from
differences of sections extending into that filling.
\end{theorem}

\begin{proof}
We distinguish an individual affine origin from its
radial-loop class throughout. B.4.29 prevents an
argument based only on extending invariant forms
from controlling an individual origin. B.4.32 is
the additional global input that controls the loop
class without making such an assertion.

\emph{Reduction of an arbitrary filling to a marked proper map.}
Suppose a filling exists. Its punctured part is a
connected smooth fourfold. Retain its unique reduced
main component, normalize, and resolve it properly,
with every centre over the central parameter. The
result is a smooth connected fourfold proper over
the parameter disc and unchanged over its puncture.
The parameter is a nonzero divisor, so its map to
the disc is flat. The complex-analytic resolution
used here has the scope recorded in B.4.31 and
\citet[Theorems 1.6 and 1.10]{BM95AB}.

B.4.4 extends the intrinsically marked map to
$\PP^1$ meromorphically. Resolve its proper graph
and the support of the central divisor. We obtain
a proper holomorphic map
\[
 (f,r):Z\longrightarrow\PP^1\times\Delta_r
 \tag{B.4.33.1}
\]
with $Z$ smooth, whose restriction over $r\ne0$
is exactly the marked actual family. The central
support is simple normal crossing after further
embedded resolution, but its multiplicities are
retained. B.4.31 and B.4.32 apply in this scope.
There is no use of a semistable-reduction theorem
for threefolds.

A hypothetical filling remains proper after any
further finite base change, normalization, and
retention of its main component. A finite map of
disc germs is a coordinate power after absorbing
a holomorphic unit into the source coordinate.
We may therefore make a further cover so that
the actual parameter is $w=r^6$, followed, if
necessary, by $r=R^n$ with $n>0$. This notation
includes every original positive cover degree:
pulling $w=\rho^N$ back by $\rho=R^6$ gives
$w=R^{6N}$ and hence $r=R^N$.

Put $q=1/f$ near the marked cusp and make the
normalized base change $q=Q^{12m}$, $m>0$.
Write
\[
 B=\Delta_Q\times\Delta_R,\qquad
 B^\times=\Delta_Q^*\times\Delta_R^*,\qquad
 B_q=\Delta_Q^*\times\Delta_R,\quad
 B_r=\Delta_Q\times\Delta_R^*.
 \tag{B.4.33.2}
\]
The proper normalized main component above $B$
has the actual torus torsor $\mathscr P$ on
$B^\times$. Its Jacobian is the restriction of
the specific log torus $\mathcal A_B$ of
B.4.24. The completion on $B_q$ is unknown;
we do not identify it with the rotating model.
On $B_r$ it is the prescribed actual cusp,
up to the proper modifications just made.

\emph{Only the origin can have excess fibre dimension.}
There are finitely many irreducible components
of the compact central divisor in (B.4.33.1).
On each component the marked map is either
constant or has two-dimensional level sets.
Thus a three-dimensional fibre of $(f,r)$
can occur only at one of the finitely many
constant images. Shrink the cusp disc to
exclude all these images except its centre.
Also exclude the finitely many noncusp values
removed in B.4.31--B.4.32. The map is now
equidimensional of relative dimension two
on $B\setminus\{0\}$. Finite base change
and normalization preserve this assertion.

Apply the coefficient-ideal construction of
B.4.30 to its possible three-dimensional
origin fibre. It supplies a finite sequence
of point blowups
\[
 \beta:S\longrightarrow B
 \tag{B.4.33.3}
\]
with centres over the origin, and a proper
normalized main transform whose every fibre
has dimension two. We use precisely the
convergent finite-projection argument proved
there, so this is not an invocation of a
projective flattening theorem. Put
$K=\beta^{-1}(0)$. When no blowup is needed,
$K$ is a point; otherwise it is the compact
exceptional tree. Write $E$ for the exceptional
divisor, possibly empty, and $D_Q,D_R$ for
the strict transforms of the axes.

At every point of $K$, a transverse analytic
slice gives a finite local multisection.
Normalize it; a normal surface is
Cohen--Macaulay, so its finite map to the
smooth surface base is flat. Choose a finite
open cover $U_i$ of a neighbourhood of $K$
and such multisections of degrees $d_i>0$.
Properness of $\beta$ implies that this
neighbourhood contains the inverse image
of a smaller bidisc. After shrinking, the
$U_i$ cover all of $S$. Consequently all
degree, resolution, and norm choices below
can be made from one finite collection.

Let $d$ be a common multiple of the $d_i$.
On the ordinary locus of $U_i$, the trace
of its multisection multiplied by $d/d_i$
is a section $t_i$ of the contracted
$d$-fold torsor $\mathscr P^d$. A trace
uses the group law of the Jacobian and
does not require an affine origin. Thus
$t_i-t_j$ is an ordinary Jacobian section
on every overlap.

\emph{Extending the differences at every boundary divisor.}
Let $\mathcal A_S=\beta^*\mathcal A_B$ in
the associated-log-torus sense of B.4.27.
Its exact bounded-character sheaf is
\[
 \mathcal T_{\log}^{(Y)}
  =\mathcal M_{D_Q+E}^{\rm gp}
        \times\mathcal M_{D_Q+D_R+E}^{\rm gp},
 \qquad Y=\Z\widehat u\oplus\Z\widehat\gamma.
 \tag{B.4.33.4}
\]
We verify this specific target for the
extended differences, rather than assigning
a log-group structure to the unknown cap.

At a generic exceptional component, the
orders of both $Q$ and $R$ are positive.
The actual monodromy has full rank two;
its primitive invariant lattice is
$\Z\widehat w\oplus\Z\widehat\delta$.
The saturation, integral-residue, and
convergent-primitive proof of B.4.30
therefore extends differences of two
extending sections to $\mathcal A_S$.
Every integer character order is bounded
at such a divisor. The same full-rank
argument applies at a generic point of
$D_Q$; here the underlying punctured-axis
family is additionally the actual given
cusp model.

At a generic point of $D_R$, the monodromy
has rank one. The saturation step of
B.4.30 still applies to the actual
rank-two subbundle of normalized invariant
one-forms:
the logarithmic term is solely
$\widehat\delta\otimes\gamma$, and the
$\widehat w,\widehat\delta$ columns
normalize a holomorphic subbundle of the
canonical log-de Rham extension. In this
case we do not replace its three-dimensional
monodromy kernel by a two-dimensional one.
Instead B.4.31 supplies the needed stronger
geometric fact: every local angular
meridian of the cap is an integral multiple
of $\widehat\delta$. After any local
semistable cover and toric subdivision,
the same holds for the new angular
meridians: their rational images are
combinations of the previous angular
images, and the primitive cycle
$\widehat\delta$ makes an integral image
an integral multiple of it.

Consequently the first normalized form
has zero local residues, and the second
has integral residues. The explicit
relative primitives of B.4.30 exponentiate
to
\[
             \mathcal O^*\times
                    \mathcal M_{D_R}^{\rm gp}.
 \tag{B.4.33.5}
\]
This is exactly the restriction of
(B.4.33.4): the first character has zero
order at the strict $R$ axis. The finite
chart-chain comparison works because the
proper central surface fibres are
connected. Indeed the normal proper map
to the normal surface base has trivial
finite Stein factor, since its ordinary
fibres are connected tori. Normalization
and retention of the main component after
a root have the same property.

These statements concern local reductions
near a generic point of a boundary curve.
B.4.30 gives them holomorphically in the
transverse parameter, with a finite
common root before the toric subdivision.
The root simultaneously accounts for
central multiplicities and splitting of
the chosen finite multisections. Norms
then show that one positive multiple of
each trace difference extends downstairs.
Its first order at $D_R$ remains zero:
the norm of a holomorphic unit is a unit.
Choose one common positive multiplier
for all the finitely many curves and
traces, and replace $d,t_i$ accordingly.
Denote their final common degree by $D>0$.
For precision, resolve the inverse image
of the full boundary over a relatively
compact neighbourhood of $K$; properness
then supplies finitely many divisor
multiplicities and finitely many trace
degrees. Their multiplicity common multiple
times $\operatorname{lcm}(1,\ldots,d_{\max})$
works for every local transverse root;
no global root along a boundary curve
is required, because norms are taken
locally and the unique extended sections
glue. Other branch curves of the finite
multisections meet a boundary component
only at isolated points unless they
equal it; after shrinking around the
compact set these intersections are
finite. The split sections lift through
proper modifications meromorphically from
their smooth surface bases, with only
codimension-two indeterminacy, and these
isolated points are treated by the
Hartogs argument below.

The only remaining exceptional points of
the boundary are isolated. The Hartogs
argument in B.4.30 applies to sections
of exactly the sheaf (B.4.33.4). On a
punctured small bidisc the constant
period sheaf $Y$ has zero $H^1$, so a
section lifts to bounded characters.
After division by their boundary monomials,
both units and inverses extend by ordinary
Hartogs. Restoring the monomials preserves
the zero first-character order on $D_R$.
The extension is unique modulo $Y$.
Thus every $t_i-t_j$ extends to an
$\mathcal A_S$-section on $U_i\cap U_j$,
and uniqueness preserves the cocycle
identities.

\emph{The cap-derived torsor has a global origin.}
Glue the trivial $\mathcal A_S$-torsors
on $U_i$ by these extended differences.
The result is a torsor $\mathcal Q$ on
\emph{all} of $S$. Its ordinary restriction
is canonically $\mathscr P^D$, using
the actual trace sections as the local
identifications. It has not been required
to match the prescribed rotating model
on $D_R$. B.4.27 gives
\[
 H^1(S,\mathcal A_S)=0.
 \tag{B.4.33.6}
\]
Choose a global section $h$ of $\mathcal Q$.
On the ordinary locus it is a section
of the actual $D$-fold affine torsor.
Near a generic point of $D_R$ its
difference from a multiplied local cap
trace is an $\mathcal A_S$-section.

\emph{The strict-$R$ comparison controls only the loop class.}
For an ordinary Jacobian section $a$
write $\nu_R(a)$ for its radial connecting
class in
\[
 V=(\Lambda\otimes\Q)/\Q\widehat\delta,
 \tag{B.4.33.7}
\]
as defined in B.4.32. Let $\ell_\gamma$
be the integral covector taking value one
on $\widehat\gamma$ and zero on
$\widehat u,\widehat w,\widehat\delta$.
It descends to $V$ and is invariant under
the actual monodromy.

Let $o_{\rm ord}$ be a B.4.19 ordinary
local reference origin. B.4.32, including
its degree-normalized trace and root
formula, gives
\[
 \nu_R(t_i-D o_{\rm ord})=0.
 \tag{B.4.33.8}
\]
An $\mathcal A_S$-section near $D_R$
has a local lift (B.4.33.5). Its radial
loop can therefore add only a multiple
of $\widehat\delta$. It follows that
$\nu_R(h-D o_{\rm ord})=0$ as well.

The rotating origin $o$ of B.4.24 on
$B_q$ has the same radial-loop class
as $o_{\rm ord}$. This is the exact
comparison (B.4.19.5): their coordinate
change is $G(s)+6v\widehat\delta$,
where $G$ is independent of $r$.
After $r=R^n$ its radial defect is still
purely $\widehat\delta$. Therefore
\[
 \ell_\gamma\nu_R(h-D o)=0.
 \tag{B.4.33.9}
\]
No boundedness or meromorphic extension
of $h-D o$ is asserted or needed.

\emph{The actual cusp origin gives the second zero.}
The section $\sigma$ of (B.4.24.2)
is an actual extending origin on $B_r$.
Its characters are the actual unit
characters of the known cusp fan.
Over $B_r$ the cap still projects properly
to that prescribed fan presentation,
so comparing a finite trace with
$d_i\sigma$ is exactly the character-norm
calculation of B.4.26. This does not
require a global lift of $\sigma$
through every modification; alternatively
such a lift is meromorphic and is
holomorphic away from isolated points,
where the same Hartogs argument applies.
Include these bounded-degree norm
multipliers in the choice of $D$ above.
Hence $D\sigma$ is a
section of $\mathcal Q|_{B_r}$ and
\[
                  h-D\sigma\in
                         \mathcal A_B(B_r).
 \tag{B.4.33.10}
\]
Here $S$ is identified with $B$ away
from its origin, so there is no base
modification in this restriction.

The $Q$-ghost map on the period sheaf
along the cusp is
\[
 Y\longrightarrow\Z^2,\qquad
 \widehat u\longmapsto(12m,0),\quad
 \widehat\gamma\longmapsto(0,12m).
 \tag{B.4.33.11}
\]
It is injective. The ghost-cocycle
argument of B.4.24 therefore makes the
radial $Y$-connecting class of every
$\mathcal A_B(B_r)$-section zero.
In particular
\[
             \ell_\gamma\nu_R(h-D\sigma)=0.
 \tag{B.4.33.12}
\]
The quotient by the cocharacter cycles
$\widehat w,\widehat\delta$ in that
argument does not change this evaluation.

\emph{The exact actual clutching gives a contradiction.}
On $B^\times$, the coordinates based at
$\sigma$ and $o$ differ by the exact
translation $F$ of (B.4.24.3):
\[
 \zeta\longmapsto\zeta-\sigma+F.
\]
Thus $\sigma-o=F$ in the common Jacobian,
and
\[
 (h-D o)-(h-D\sigma)=D F.
 \tag{B.4.33.13}
\]
With $s=(2\pi i)^{-1}\log q$ and
$v=(2\pi i)^{-1}\log r$, the exact
period difference is
\[
 F(s,v+n)-F(s,v)
       =n\gamma-6n(n-1)\delta.
 \tag{B.4.33.14}
\]
The corresponding integral period label
has $\ell_\gamma$ value $n$.
Applying $\ell_\gamma\nu_R$ to
(B.4.33.13), equations (B.4.33.9) and
(B.4.33.12) give zero on the left,
whereas (B.4.33.14) gives $Dn>0$ on
the right. This is impossible.

All auxiliary coordinate roots, toric
norms, and trace multipliers have positive
degree and merely multiply this nonzero
integer. The contradiction excludes the
initial arbitrary proper analytic filling.
In particular the smooth reduced-NC
filling requested in Part B does not
exist for any positive cover exponent.
\end{proof}

The final step uses two separate facts: a
cap-derived origin has zero global radial
defect by B.4.32, and a section extending
along the actual cusp has zero radial
ghost class by B.4.24. Neither fact
identifies the unknown central model with
B.4.19. A common essential translation
with zero loop class, as allowed by
B.4.29, is harmless in this comparison.
The untwisted Engel family provides a
useful scope check: its identity closing
has zero affine origins on both axial
fan models, hence $F=0$. The nonzero
clutching class used in (B.4.33.14)
is specific to the actual twisted
family; the proof does not exclude
the untwisted family's known filling.
